\chapter{Software Architecture and HPC Optimizations}\label{chapter:software-arch-and-hpc-optimizations}



\section{Bitwise operations for Morton Encoding} \label{sec:bitwise-operations-for-morton-encoding}
As detailed in Section \ref{subsubsec:z-order-curves}, the Morton Z-order curve is used to map the 3D coordinates of the mesh elements for a more effient octree construction. To optimize this process without relying on expensive looping, we use a series of bitwise operations that interleave the bits of the binary representations of the $x$, $y$, and $z$ coordinates.

The basic algorithm that we use to compute the Morton key for a given 3D coordinate is as follows:
\begin{enumerate} [label=\textbf{Step \arabic*.}, leftmargin=*]
    \item \textbf{Grid Quantization:} We define a discrete grid across the global bounding box. With a 64 bit integer, we can divide each spatial dimension into $2^{20} \left(1,048,576\right)$ bins, which allows us to represent up to $2^{60}$ unique spatial locations, more than enough for our problem size.
    \item \textbf{Coordinate Normalization:} We normalize the $x$, $y$, and $z$ coordinates of each mesh element to fit within $[0, 1]$ relative to the global bounding box, and then multiply by the number of bins to get the quantized integer coordinates.
    \item \textbf{Bit Truncation:} We truncate the quantized coordinates to fit within 21 bits, ensuring that they can be interleaved into a single 64-bit integer without overflow. This is achieved by applying a bitwise AND operation with a mask that retains only the lower 21 bits of each coordinate.
    \item \textbf{Bit Expansion:} We use a sequence of bitwise shifts and bitwise operations with bit masks to add two zero bits between each bit of the quantized coordinates. This process is repeated for each coordinate.
    \item \textbf{Bit Interleaving:} Finally, we combine the expanded bits of each coordinate by shifting the expanded $y$ bits by one, and the expanded $z$ bits by two, and then performing a bitwise OR operation to interleave the bits together, resulting in the final Morton key for that mesh element.
\end{enumerate} 
The basics algorithm of this process is shown for the 2D case in \cite{bittwiddlinghacks}
\section{Caching the Matrices} \label{sec:caching-the-matrices}

\lipsum[1-1]

\section{OpenMP speed up} \label{sec:openmp-speed-up}

\lipsum[1-1]

\section{Solving the system - vmult + GMRES + The Real-Equivalent system} \label{sec:solving-the-system}

\lipsum[1-1]